% ==== SR: границы a(n) по режимам, n = 2..20 ====
\def\vSRACons{(2,3.000000) (3,6.000000) (4,10.000000) (5,15.000000) (6,21.000000) (7,28.000000) (8,36.000000) (9,45.000000) (10,55.000000) (11,64.000000) (12,72.000000) (13,81.000000) (14,91.000000) (15,100.000000) (16,110.000000) (17,121.000000) (18,133.000000) (19,144.000000) (20,156.000000)}
\def\vSRASym{(2,3.000000) (3,6.000000) (4,10.000000) (5,15.000000) (6,21.000000) (7,28.000000) (8,36.000000) (9,45.000000) (10,55.000000) (11,66.000000) (12,78.000000) (13,87.000000) (14,97.000000) (15,106.000000) (16,118.000000) (17,127.000000) (18,139.000000) (19,150.000000) (20,162.000000)}
\def\vSRANear{(2,2.000000) (3,5.000000) (4,9.000000) (5,14.000000) (6,18.000000) (7,23.000000) (8,27.000000) (9,32.000000) (10,38.000000) (11,43.000000) (12,49.000000) (13,56.000000) (14,62.000000) (15,69.000000) (16,75.000000) (17,82.000000) (18,90.000000) (19,97.000000) (20,105.000000)}
\def\vSRANpExact{(2,3.000000) (3,6.000000) (4,10.000000) (5,15.000000) (6,21.000000) (7,28.000000) (8,36.000000) (9,45.000000) (10,53.000000) (11,62.000000) (12,70.000000) (13,79.000000) (14,89.000000) (15,98.000000) (16,108.000000) (17,119.000000) (18,131.000000) (19,142.000000) (20,154.000000)}
\def\vSRANpAsym{(2,6.708152) (3,11.224886) (4,16.431551) (5,22.248425) (6,28.617957) (7,35.496206) (8,42.848242) (9,50.645441) (10,58.863799) (11,67.482814) (12,76.484706) (13,85.853870) (14,95.576461) (15,105.640091) (16,116.033588) (17,126.746809) (18,137.770493) (19,149.096140) (20,160.715910)}
% ==== SR: достигнутый FAR(n) по режимам, n = 2..20 ====
\def\vSRFARCons{(2,0.500000) (3,0.250000) (4,0.125000) (5,0.062500) (6,0.031250) (7,0.015625) (8,0.007812) (9,0.003906) (10,0.001953) (11,0.001953) (12,0.002441) (13,0.002441) (14,0.002319) (15,0.002625) (16,0.002686) (17,0.002579) (18,0.002335) (19,0.002399) (20,0.002325)}
\def\vSRFARSym{(2,0.500000) (3,0.250000) (4,0.125000) (5,0.062500) (6,0.031250) (7,0.015625) (8,0.007812) (9,0.003906) (10,0.001953) (11,0.000977) (12,0.000488) (13,0.000732) (14,0.000854) (15,0.001160) (16,0.001007) (17,0.001343) (18,0.001289) (19,0.001411) (20,0.001432)}
\def\vSRFARNear{(2,0.500000) (3,0.250000) (4,0.125000) (5,0.062500) (6,0.062500) (7,0.046875) (8,0.054688) (9,0.054688) (10,0.048828) (11,0.053711) (12,0.052246) (13,0.047852) (14,0.049438) (15,0.047913) (16,0.050659) (17,0.050537) (18,0.048279) (19,0.049366) (20,0.048441)}
% ==== SR: таблица режимов, n = 2..20 ====
\def\vSRModeTable{%
2 & 3 & 0.500000 & 3 & 0.500000 & 2 & 0.500000 \\
3 & 6 & 0.250000 & 6 & 0.250000 & 5 & 0.250000 \\
4 & 10 & 0.125000 & 10 & 0.125000 & 9 & 0.125000 \\
5 & 15 & 0.062500 & 15 & 0.062500 & 14 & 0.062500 \\
6 & 21 & 0.031250 & 21 & 0.031250 & 18 & 0.062500 \\
7 & 28 & 0.015625 & 28 & 0.015625 & 23 & 0.046875 \\
8 & 36 & 0.007812 & 36 & 0.007812 & 27 & 0.054688 \\
9 & 45 & 0.003906 & 45 & 0.003906 & 32 & 0.054688 \\
10 & 55 & 0.001953 & 55 & 0.001953 & 38 & 0.048828 \\
11 & 64 & 0.001953 & 66 & 0.000977 & 43 & 0.053711 \\
12 & 72 & 0.002441 & 78 & 0.000488 & 49 & 0.052246 \\
13 & 81 & 0.002441 & 87 & 0.000732 & 56 & 0.047852 \\
14 & 91 & 0.002319 & 97 & 0.000854 & 62 & 0.049438 \\
15 & 100 & 0.002625 & 106 & 0.001160 & 69 & 0.047913 \\
16 & 110 & 0.002686 & 118 & 0.001007 & 75 & 0.050659 \\
17 & 121 & 0.002579 & 127 & 0.001343 & 82 & 0.050537 \\
18 & 133 & 0.002335 & 139 & 0.001289 & 90 & 0.048279 \\
19 & 144 & 0.002399 & 150 & 0.001411 & 97 & 0.049366 \\
20 & 156 & 0.002325 & 162 & 0.001432 & 105 & 0.048441 \\}
% ==== SR: пример n = 14 (критерий nearest), a = 40..80 ====
\def\vSRNearAbs14{(40,0.216553) (41,0.216553) (42,0.193726) (43,0.193726) (44,0.172607) (45,0.172607) (46,0.153076) (47,0.153076) (48,0.135254) (49,0.135254) (50,0.118896) (51,0.118896) (52,0.104004) (53,0.104004) (54,0.090576) (55,0.090576) (56,0.078491) (57,0.078491) (58,0.067627) (59,0.067627) (60,0.057983) (61,0.057983) (62,0.049438) (63,0.049438) (64,0.041870) (65,0.041870) (66,0.035278) (67,0.035278) (68,0.029541) (69,0.029541) (70,0.024536) (71,0.024536) (72,0.020264) (73,0.020264) (74,0.016602) (75,0.016602) (76,0.013428) (77,0.013428) (78,0.010742) (79,0.010742) (80,0.008545)}
\def\vSRNearMin14{(62,0.049438) (63,0.049438)}
% ==== SR: мощность при сдвиге (n = 16), delta = 0..2 ====
% n=16: книга a=110, npSPC exact a=108
\def\vSROCBk{(0.0,0.002686) (0.25,0.014925) (0.5,0.086417) (0.75,0.295649) (1.0,0.623783) (1.25,0.882588) (1.5,0.980478) (1.75,0.998228) (2.0,0.999897)}
\def\vSROCNp{(0.0,0.003357) (0.25,0.017994) (0.5,0.100213) (0.75,0.328371) (1.0,0.662585) (1.25,0.904110) (1.5,0.985944) (1.75,0.998910) (2.0,0.999948)}
% ==== SR: распределение SR в IC (n = 5) ====
\def\vSRpmfFive{(-15,0.031250) (-13,0.031250) (-11,0.031250) (-9,0.062500) (-7,0.062500) (-5,0.093750) (-3,0.093750) (-1,0.093750) (1,0.093750) (3,0.093750) (5,0.093750) (7,0.062500) (9,0.062500) (11,0.031250) (13,0.031250) (15,0.031250)}
% ==== SR: односторонний хвост P(SR>=a) для n = 5,10,14,20 ====
\def\vSRTailFive{(15,0.031250) (13,0.062500) (11,0.093750) (9,0.156250) (7,0.218750) (5,0.312500) (3,0.406250) (1,0.500000)}
\def\vSRTailTen{(55,0.000977) (53,0.001953) (51,0.002930) (49,0.004883) (47,0.006836) (45,0.009766) (43,0.013672) (41,0.018555) (39,0.024414) (37,0.032227) (35,0.041992) (33,0.052734) (31,0.065430) (29,0.080078) (27,0.096680) (25,0.116211) (23,0.137695) (21,0.161133) (19,0.187500) (17,0.215820) (15,0.246094) (13,0.278320) (11,0.312500) (9,0.347656) (7,0.384766) (5,0.422852) (3,0.460938) (1,0.500000)}
\def\vSRTailFourteen{(105,0.000061) (103,0.000122) (101,0.000183) (99,0.000305) (97,0.000427) (95,0.000610) (93,0.000854) (91,0.001160) (89,0.001526) (87,0.002014) (85,0.002625) (83,0.003357) (81,0.004272) (79,0.005371) (77,0.006714) (75,0.008301) (73,0.010132) (71,0.012268) (69,0.014771) (67,0.017639) (65,0.020935) (63,0.024719) (61,0.028992) (59,0.033813) (57,0.039246) (55,0.045288) (53,0.052002) (51,0.059448) (49,0.067627) (47,0.076538) (45,0.086304) (43,0.096863) (41,0.108276) (39,0.120605) (37,0.133789) (35,0.147888) (33,0.162903) (31,0.178772) (29,0.195496) (27,0.213135) (25,0.231567) (23,0.250793) (21,0.270813) (19,0.291504) (17,0.312866) (15,0.334900) (13,0.357422) (11,0.380432) (9,0.403870) (7,0.427612) (5,0.451599) (3,0.475769) (1,0.500000)}
\def\vSRTailTwenty{(209,0.000001) (207,0.000002) (205,0.000003) (203,0.000005) (201,0.000007) (199,0.000010) (197,0.000013) (195,0.000018) (193,0.000024) (191,0.000031) (189,0.000041) (187,0.000052) (185,0.000067) (183,0.000084) (181,0.000105) (179,0.000131) (177,0.000161) (175,0.000197) (173,0.000241) (171,0.000293) (169,0.000354) (167,0.000425) (165,0.000508) (163,0.000605) (161,0.000716) (159,0.000845) (157,0.000993) (155,0.001163) (153,0.001356) (151,0.001576) (149,0.001827) (147,0.002110) (145,0.002430) (143,0.002790) (141,0.003195) (139,0.003648) (137,0.004154) (135,0.004718) (133,0.005344) (131,0.006040) (129,0.006808) (127,0.007656) (125,0.008591) (123,0.009617) (121,0.010742) (119,0.011975) (117,0.013321) (115,0.014788) (113,0.016384) (111,0.018117) (109,0.019995) (107,0.022027) (105,0.024220) (103,0.026585) (101,0.029129) (99,0.031861) (97,0.034790) (95,0.037926) (93,0.041275) (91,0.044847) (89,0.048654) (87,0.052699) (85,0.056993) (83,0.061546) (81,0.066363) (79,0.071453) (77,0.076823) (75,0.082479) (73,0.088427) (71,0.094674) (69,0.101225) (67,0.108084) (65,0.115256) (63,0.122744) (61,0.130549) (59,0.138678) (57,0.147126) (55,0.155897) (53,0.164991) (51,0.174405) (49,0.184138) (47,0.194188) (45,0.204549) (43,0.215217) (41,0.226188) (39,0.237453) (37,0.249004) (35,0.260837) (33,0.272938) (31,0.285298) (29,0.297910) (27,0.310757) (25,0.323828) (23,0.337111) (21,0.350591) (19,0.364253) (17,0.378083) (15,0.392063) (13,0.406178) (11,0.420411) (9,0.434744) (7,0.449159) (5,0.463639) (3,0.478165) (1,0.492718)}
